% Source for the figure of LinkCutTree (sec. 8.13): a tree of 10 nodes
% rooted at 0 after access(6) and access(9), its preferred paths in violet,
% then access(7), the edges it makes preferred in orange, and the splay
% trees of the paths after it, the path-parent pointers dashed; the state
% printed by a program on the class. Build: make figures (in ../build).
\documentclass[border=4pt]{standalone}
\usepackage[T1]{fontenc}
\usepackage{lmodern}
\usepackage{amsmath,amssymb}
\usepackage{tikz}
\usepackage{lamalunga}
% tikz-ads.tex: the TikZ styles shared by the figures of the memory of the
% manual, \input by each standalone figure: a cell of a vector (cell), the
% cell that an operation touches (touched), a free cell (free), a node of a
% linked structure (node), the labels of indices (index) and of variables
% (var), a pointer (ptr), and \cells and \indices to draw a row of cells.
\usetikzlibrary{arrows.meta,positioning,calc,fit,shapes.multipart}
\definecolor{adsTouched}{RGB}{214,196,236}
\definecolor{adsFree}{RGB}{236,236,236}
\tikzset{
  ads/.style={>={Stealth[length=2mm]},font=\small,
              every node/.append style={inner sep=0pt}},
  cell/.style={draw,minimum width=7mm,minimum height=7mm,
               font=\small\ttfamily,outer sep=0pt},
  touched/.style={cell,fill=adsTouched},
  free/.style={cell,fill=adsFree,text=black!45},
  node/.style={draw,rounded corners=1pt,minimum width=8mm,
               minimum height=6mm,font=\small\ttfamily},
  index/.style={font=\scriptsize\ttfamily,text=black!60},
  var/.style={font=\small\itshape},
  ptr/.style={->,thick}
}
% \cells{name}{x,y}{list}: a row of cells, the first at (x,y), named
% name-0, name-1, ...; each item is text/style, the style one of cell,
% touched and free (e.g., {a/cell,e/touched,/free})
\newcommand{\cells}[3]{%
  \foreach \v/\s [count=\i from 0] in {#3}
    \node[\s] (#1-\i) at ($(#2)+(\i*7mm,0)$) {\v};}
% \indices{name}{n}: the indices 0 .. n-1 under the cells of name
\newcommand{\indices}[2]{%
  \foreach \i in {0,...,\numexpr#2-1\relax}
    \node[index,below=1.2mm of #1-\i] {\i};}

\begin{document}
\begin{tikzpicture}[ads]
  % (a) the tree before, preferred paths 0-2-5-9 and 1-3-6
  \begin{scope}
    \node[vertex] (a0) at (30mm,0) {0};
    \node[vertex] (a1) at (15mm,-13mm) {1};
    \node[vertex] (a2) at (45mm,-13mm) {2};
    \node[vertex] (a3) at (8mm,-26mm) {3};
    \node[vertex] (a4) at (22mm,-26mm) {4};
    \node[vertex] (a5) at (45mm,-26mm) {5};
    \node[vertex] (a6) at (8mm,-39mm) {6};
    \node[vertex] (a7) at (22mm,-39mm) {7};
    \node[vertex] (a8) at (38mm,-39mm) {8};
    \node[vertex] (a9) at (52mm,-39mm) {9};
    \draw[edge] (a0) -- (a1);
    \draw[edge-touched] (a0) -- (a2);
    \draw[edge-touched] (a1) -- (a3);
    \draw[edge] (a1) -- (a4);
    \draw[edge-touched] (a2) -- (a5);
    \draw[edge-touched] (a3) -- (a6);
    \draw[edge] (a4) -- (a7);
    \draw[edge] (a5) -- (a8);
    \draw[edge-touched] (a5) -- (a9);
    \node[var] at (30mm,-50mm) {(a) before};
  \end{scope}
  % (b) the tree after access(7)
  \begin{scope}[xshift=70mm]
    \node[vertex] (b0) at (30mm,0) {0};
    \node[vertex] (b1) at (15mm,-13mm) {1};
    \node[vertex] (b2) at (45mm,-13mm) {2};
    \node[vertex] (b3) at (8mm,-26mm) {3};
    \node[vertex] (b4) at (22mm,-26mm) {4};
    \node[vertex] (b5) at (45mm,-26mm) {5};
    \node[vertex] (b6) at (8mm,-39mm) {6};
    \node[vertex-touched] (b7) at (22mm,-39mm) {7};
    \node[vertex] (b8) at (38mm,-39mm) {8};
    \node[vertex] (b9) at (52mm,-39mm) {9};
    \draw[edge-added] (b0) -- (b1);
    \draw[edge] (b0) -- (b2);
    \draw[edge] (b1) -- (b3);
    \draw[edge-added] (b1) -- (b4);
    \draw[edge-touched] (b2) -- (b5);
    \draw[edge-touched] (b3) -- (b6);
    \draw[edge-added] (b4) -- (b7);
    \draw[edge] (b5) -- (b8);
    \draw[edge-touched] (b5) -- (b9);
    \node[var] at (30mm,-50mm) {(b) after \texttt{access(7)}};
  \end{scope}
  % (c) the splay trees after access(7) and the path-parent pointers
  \begin{scope}[xshift=40mm,yshift=-62mm]
    \node[vertex-touched] (c7) at (40mm,0) {7};
    \node[vertex] (c0) at (16mm,-11mm) {0};
    \node[vertex] (c4) at (32mm,-22mm) {4};
    \node[vertex] (c1) at (24mm,-33mm) {1};
    \node[vertex] (c5) at (2mm,-24mm) {5};
    \node[vertex] (c2) at (-6mm,-35mm) {2};
    \node[vertex] (c9) at (10mm,-35mm) {9};
    \node[vertex] (c8) at (2mm,-48mm) {8};
    \node[vertex] (c3) at (24mm,-46mm) {3};
    \node[vertex] (c6) at (32mm,-57mm) {6};
    \draw[edge-added] (c7) -- (c0);
    \draw[edge-added] (c0) -- (c4);
    \draw[edge-added] (c4) -- (c1);
    \draw[edge-touched] (c5) -- (c2);
    \draw[edge-touched] (c5) -- (c9);
    \draw[edge-touched] (c3) -- (c6);
    \draw[guide,->] (c5) -- (c0);
    \draw[guide,->] (c3) -- (c1);
    \draw[guide,->] (c8) -- (c5);
    \node[var,anchor=west] at (46mm,-24mm) {(c) the splay trees};
  \end{scope}
\end{tikzpicture}
\end{document}
